\documentclass[11pt]{article}
\usepackage[T1]{fontenc}
\usepackage[utf8]{inputenc}
\usepackage{lmodern,microtype}
\usepackage[letterpaper,margin=0.88in]{geometry}
\usepackage{amsmath,amssymb,amsthm,mathtools}
\usepackage{booktabs,longtable,array,enumitem}
\usepackage[dvipsnames]{xcolor}
\usepackage{fancyhdr}
\usepackage[colorlinks=true,linkcolor=MidnightBlue,urlcolor=MidnightBlue,citecolor=MidnightBlue,breaklinks=true]{hyperref}
\usepackage{listings}
\hypersetup{pdftitle={Shared Binomial Certificates and Short Carriers for Diophantine Consistency},pdfauthor={Research report prepared for Vladimir Reshetnikov}}
\newcommand{\N}{\mathbb N}
\newcommand{\Z}{\mathbb Z}
\newcommand{\Q}{\mathbb Q}
\newcommand{\Con}{\operatorname{Con}}
\newcommand{\ZFC}{\mathsf{ZFC}}
\newcommand{\code}[1]{\texttt{\small #1}}
\newtheorem{theorem}{Theorem}[section]
\newtheorem{lemma}[theorem]{Lemma}
\newtheorem{proposition}[theorem]{Proposition}
\newtheorem{corollary}[theorem]{Corollary}
\theoremstyle{definition}\newtheorem{definition}[theorem]{Definition}
\theoremstyle{remark}\newtheorem{remark}[theorem]{Remark}
\setlength{\parskip}{4pt}
\setlength{\parindent}{0pt}
\setlength{\emergencystretch}{2em}
\setlist{topsep=4pt,itemsep=3pt,parsep=0pt}
\pagestyle{fancy}\fancyhf{}
\fancyhead[L]{\small Diophantine consistency}
\fancyhead[R]{\small 24 October 2026}
\fancyfoot[C]{\thepage}
\setlength{\headheight}{14pt}
\lstset{basicstyle=\ttfamily\small,breaklines=true,columns=fullflexible,keepspaces=true,showstringspaces=false,frame=single,rulecolor=\color{black!20}}
\begin{document}
\hypersetup{pageanchor=false}
\begin{titlepage}
\vspace*{0.7in}
{\color{MidnightBlue}\large PROVEIT RESEARCH REPORT\par}
\vspace{0.35in}
{\Huge\bfseries Shared Binomial Certificates\par}
\vspace{0.12in}
{\Huge\bfseries and Short Carriers\par}
\vspace{0.3in}
{\LARGE Explicit Diophantine forms of $\Con(\ZFC)$\par}
\vspace{0.45in}
{\large Prepared for Vladimir Reshetnikov\par}
{\large 24 October 2026\par}
\vspace{0.5in}
\begin{minipage}{0.94\linewidth}
\textbf{Main outputs.} A direct quartic in 319 nonnegative integers; an
eleven-integer equation of degree 966 with a 111-operation evaluator; and
an eleven-integer equation of degree 440 with a 126-operation evaluator.
Every coefficient and arithmetic operation is supplied in the companion
archive. The new constructions preserve the finite-refutation endpoint of
the explicitly retained ordinary source proofs.

\medskip
\textbf{Main idea.} Reuse the same binomial certificate to certify a
factorial and the product of the CRT moduli. On the small-variable side,
move a coprime square condition into a Pell coefficient, and use a short
carrier whose weaker divisibility is sufficient for exact recovery.

\medskip
\textbf{Scope.} These are concrete improvements against the inspected
repository and all eight incoming archives. No global minimum, historical
priority theorem, proof of consistency, or newly completed Lean/Rocq
endpoint is claimed. The arithmetic proofs, finite artifact checks, and
inherited semantic calibration are identified separately.
\end{minipage}
\vfill
{\small Comparison pin:\par
\texttt{45b399ec51dfbecb891dd51d8a9e4b6b7c361ac6}\par
\url{https://github.com/VladimirReshetnikov/ProveIt}}
\end{titlepage}
\hypersetup{pageanchor=true}

\begin{abstract}
We construct and export new fixed polynomial equations whose absence of
solutions is equivalent, by ordinary arguments formalizable in ZFC and the
retained source calibration, to $\Con(\ZFC)$. The direct proof-certificate
route has 319 nonnegative variables, degree four, 254 quadratic residuals,
10,283 collected monomials, and coefficient height 86,842,335. It improves
the inspected 349-variable endpoint in variables, residuals, terms, and
height. Its architecture needs only six full power certificates instead
of nineteen: a single exact binomial serves both the factorial and the
reciprocal CRT product, while common Pell residues handle eleven digit
requests. This route does not simulate K215 or specialize Jones's universal
polynomial.

A complementary native-source construction uses eleven unrestricted
integer variables. A coprime product is absorbed into the strong Pell
coefficient; congruences modulo four recover positivity of the extracted
factors. Together with a short powered-divisibility carrier this gives
degree 966 at 111 operations, or degree 882 at 112 operations. A three-root
companion gives degree 464 at 125 operations and degree 440 at 126
operations. The report supplies the mathematical definitions, both
directions of the new reductions, the exact degree arguments, the
coefficient conventions, and a reproducible audit. Existing lemmas and
unbounded source theorems are credited explicitly.
\end{abstract}
\tableofcontents
\newpage

\section{Outcome and the meaning of simplicity}
The target is a \emph{fixed} polynomial $P$ with integer coefficients and
a specified domain $D$ such that
\begin{equation}\label{eq:goal}
 \ZFC\vdash
 \bigl(\neg\exists\mathbf x\in D^m\; P(\mathbf x)=0\bigr)
 \ \longleftrightarrow\ \Con(\ZFC).
\end{equation}
Here $\N=\{0,1,2,\ldots\}$; $\Z$ means unrestricted integers. A named
polynomial intermediate is not an additional quantified variable.
The fixed coefficients are literal finite data, reconstructed without
running a refutation search. The equivalence does not assume consistency.

There is no single agreed numerical order on all such equations. Variable
count, degree, collected support, coefficient height, and arithmetic
circuit length measure different costs. An enormous integer hidden in a
coefficient is not free from an overall description-length perspective.
Accordingly we give a small set of explicitly measured tradeoffs.

\subsection{Measured improvements}
For the first table, operations count every live binary addition,
subtraction, and multiplication, including multiplication by a fixed
scalar and squaring. Reuse of an existing register is free. Fixed
coefficient leaves are reported separately.
\begin{center}\small
\begin{tabular}{@{}lrrrr@{}}\toprule
Construction over $\Z$ & Variables & Degree & Operations & Status\\\midrule
Incoming coprime fusion, direct &11&728&122&baseline\\
Incoming coprime fusion, centered&11&688&126&baseline\\
Incoming all-root fusion&11&1632&114&baseline\\\midrule
Short carrier, direct&11&488&121&this report\\
Short carrier, centered&11&464&125&this report\\
Short carrier, smaller index modulus&11&440&126&this report\\
Pell-coefficient saturation&11&966&111&this report\\
Saturation, smaller index modulus&11&882&112&this report\\\bottomrule
\end{tabular}
\end{center}
The 464/125 point improves both degree and cost relative to 688/126. The
440/126 point reduces the same baseline degree to 440 at unchanged cost. The
966/111 point improves both measures relative to 1632/114. These statements
compare the same unrestricted integer domain and the same six fixed
coefficient values.

For a collected polynomial $\sum_\alpha c_\alpha X^\alpha$, its height
is $\max_\alpha|c_\alpha|$, and its support size counts nonzero collected
coefficients. Both following equations are quartics over $\N$.
\begin{center}\small
\begin{tabular}{@{}lrrrr@{}}\toprule
Construction&Variables&Residuals&Terms&Height\\\midrule
Incoming factorial-carrier endpoint&349&268&12,768&86,915,503\\
Shared binomial and modular residues&319&254&10,283&86,842,335\\\bottomrule
\end{tabular}
\end{center}
This removes 30 variables (8.60\%), 14 residuals, and 2,485 terms
(19.46\%). The number of full power certificates drops from 19 to 6;
the number of exact binomial certificates drops from two to one. All
coefficients of both expanded quartics were independently re-expanded
from their residual lists for the comparison.

\subsection{The baseline includes the incoming archives}
The comparison uses the public repository at the pin on the title page,
including its latest proposal review, the active candidate and research
directories, and all eight ZIPs then in \code{docs/incoming}.
The newer incoming reports dominate several entries of the main comparison
table. In particular, the 349-variable quartic and the eleven-integer
688/126 and 1632/114 circuits are incoming baselines, not results of this
work. Their source paths and hashes are retained in the archive provenance.

Other architectures remain relevant. The repository already has
nine-positive-variable equations and a ten-integer construction, and it
retains a 37-integer quartic with enormous fixed coefficients. The new
319-variable quartic is therefore not a fewest-variable quartic record.
Its advantage is within a direct, explicitly expanded proof-certificate
architecture with modest numerical coefficients. The eleven-variable
native-source equations instead use six very large fixed integers.

\subsection{Attribution and research scope}
The squared-radius factorial identity used below was already proved in
package 105 and appears in the incoming synchronized-exponents work.
Common modular Pell residues occur in package 108. The reciprocal CRT
product, derived label range, and shared-factorial Horn-flow source occur
in Round 14. The new direct contribution is their complete integration:
the factorial and CRT binomials become \emph{the same certificate}, the
digit base is validated before that certificate is interpreted, and a
complete smaller quartic is exported.

On the small-variable side, powered divisibility and supported-modulus
Pell interpolation come from Round 14, and selective coprime fusion comes
from the incoming coprime-projection report. The new work combines them
with a smaller index modulus and a Pell-coefficient saturation that
recovers signs by congruences modulo four. This report makes no claim that
these ideas are absent from all mathematical literature.

\section{Logical endpoint and evidence boundaries}
The common semantic endpoint is the existence of a finite syntactic
refutation. For a fixed, effective presentation of ZFC, write
\[
 \operatorname{Ref}_{\ZFC}\equiv
 \exists p\;\operatorname{Proof}_{\ZFC}(p,\bot),\qquad
 \Con(\ZFC)\equiv\neg\operatorname{Ref}_{\ZFC}.
\]
Each construction proves $\exists P=0\leftrightarrow
\operatorname{Ref}_{\ZFC}$. Negation gives~\eqref{eq:goal}. Nothing in this
argument proves $\Con(\ZFC)$ itself or decides whether a zero exists.

\begin{theorem}[Ordinary correctness of the delivered reductions]
Compose the explicitly retained ordinary native/calculus-to-refutation
theorems with the arithmetic reductions proved in this report. The
polynomial in \code{direct/equation319/natural/equation.json} has a zero in $\N^{319}$
if and only if ZFC has a finite refutation. Each of the five selected
eleven-variable polynomials in the comparison table has a zero in
$\Z^{11}$ if and only if ZFC has a finite refutation. These finite
arithmetical arguments are formalizable in ZFC.
\end{theorem}

Three kinds of evidence must be kept separate.
\begin{enumerate}
\item \textbf{Ordinary mathematical proofs.} The report and retained
sources establish the unbounded number-theoretic and refutation
equivalences, including witness existence and all sign cases.
\item \textbf{Finite exact checks.} Independent programs authenticate
coefficient bytes, reconstruct the fixed tuple, verify the entire
polynomial expressions and elimination certificates, count operations,
and certify degrees. These checks do not decide consistency.
\item \textbf{Proof-assistant status.} ProveIt has separate formal
endpoints and extensive local proofs. This archive supplies no new
end-to-end Lean/Rocq build and no literal internal object-language ZFC
derivation of the newly displayed equations. Ordinary formalizability
is the claim in this report; a new kernel certificate is a further task.
\end{enumerate}

The direct route starts from the numerical C222 calculus and its retained
ordinary equivalence to finite ZF refutability, followed by the retained
ZF/ZFC refutation equivalence. The statement is applied before a
machine-specific simulation endpoint: the new direct builder does not
read or simulate the K215 transition table. The small-variable branch
deliberately retains the complete existing native-source coefficient
tuple and its calibration.

\section{One binomial certificate for factorial and CRT arithmetic}
\label{sec:direct319}

The strongest direct result of this investigation is a literal quartic in
319 nonnegative integers. It improves the latest inspected direct natural
equation in four measured quantities at once:
\[
\begin{array}{lrrrr}
 &\text{variables}&\text{quadratic rows}&\text{monomials}&\text{height}\\
\text{Round 14 derived-label source}&349&268&12\,768&86\,915\,503\\
\text{Merged factorial--CRT source}&319&254&10\,283&86\,842\,335.
\end{array}
\]
Both single equations have exact degree four. The comparison concerns
direct C222 equations with completely collected coefficients. It does not
assert superiority to the repository's few-variable equations with different
coefficient and degree costs.

\begin{theorem}[Explicit direct equation]\label{thm:direct319}
Let $R_1,\ldots,R_{254}\in\Z[z_0,\ldots,z_{318}]$ be the complete quadratic
residual list in the delivered file
\texttt{direct/equation319/natural/equation.json}, and put
\[
             P_{319}(z)=\sum_{j=1}^{254}R_j(z)^2.
\]
The supplied collected coefficient file defines the same polynomial.
Its degree is four, it has $10\,283$ nonzero monomials, and its height is
$86\,842\,335$. With the inherited numerical-calculus calibration,
\[
 \mathrm{ZFC}\vdash
 \bigl(\nexists z\in\N^{319}:P_{319}(z)=0\bigr)
       \ \longleftrightarrow\ \operatorname{Con}(\mathrm{ZFC}).
\]
Here $\N=\{0,1,2,\ldots\}$; all coefficients are fixed, literal integers.
\end{theorem}

The ordinary proof is given below. Its arithmetic source has 465 natural
coordinates and 400 quadratic rows; an independent exact certificate
eliminates 146 coordinates, including a nonnegativity proof for every
restoration. Expansion of all remaining squares verifies every final
coefficient. These finite checks support the literal construction; the
unbounded equivalence is proved mathematically, rather than inferred from
tests. No new end-to-end Lean proof term or emitted object-language ZFC
derivation is claimed.

\subsection{The improvement and its mathematical provenance}

The preceding reciprocal-CRT construction used $F=t!$ and
\[
 B=\binom{F^2+t}{t},
\]
but obtained the factorial through a different binomial,
$\binom{Ft^2+t}{t}$. We align these two uses of binomial arithmetic.
The squared-radius factorial fixed-point theorem was already established
in package 105, \emph{Bootstrap certificates}, Section 2 of its
\texttt{NEW\_RESULTS.md}; it is not a new theorem of this iteration.
The common Pell-residue representation was already established in package
108, \emph{Pell residues instead of power values}. The shared-factorial
source, reciprocal prime-power CRT and derived-label range come from the
latest Round 14 packet. These source proofs are retained in the archive.

The completed contribution is their integration: the \emph{same} binomial
output serves the factorial and CRT product; an independent radix bound
validates that output before either semantic interpretation is assumed;
common Pell residues eliminate eleven full numerator powers; and the
entire construction is compiled, contracted and measured. This is a new
complete equation relative to the inspected snapshot, without a claim of
worldwide priority for the ingredients.

\subsection{The inherited factorial identity, with an independent proof}

\begin{lemma}[Squared-radius factorial]\label{lem:direct-factorial}
For $t\ge6$ and positive integer $F$,
\[
 F=\left\lceil\frac{(F^2+1)^t}{\binom{F^2+t}{t}}\right\rceil
       \quad\Longleftrightarrow\quad F=t!.
\]
\end{lemma}
\begin{proof}
Put
\[
 H_t(F)=t!\prod_{i=1}^{t}\frac{F^2+1}{F^2+i}.
\]
The ceiling assumption gives $H_t(F)\le F\le t!$. If $F<t$, monotonicity
of each factor in its radius gives
\[
 H_t(F)\ge\frac{2^t}{t+1}>t>F,
\]
a contradiction. The middle inequality follows at $t=6$ and persists
by induction. Thus $F\ge t$.

Using $\prod(1-u_i)\ge1-\sum u_i$,
\[
 \prod_{i=1}^{t}\frac{F^2+1}{F^2+i}
 >1-\frac{t(t-1)}{2F^2}>\frac12.
\]
If $F<t^2$, then $H_t(F)>t!/2>t^2>F$, again impossible. Here
$t!>2t^2$ for $t\ge6$ follows immediately by induction. Hence $F\ge t^2$,
and the stronger estimate is
\[
 H_t(F)>t!\left(1-\frac1{2F}\right).
\]
If $t!\ge F+1$, the right side is at least
$(F+1)(1-1/(2F))=F+1/2-1/(2F)\ge F$, contradicting $H_t(F)\le F$.
Therefore $F=t!$. Conversely, for $F=t!$, the same estimate gives
$F-1/2<H_t(F)\le F$, so its ceiling is $F$.
\end{proof}

Package 105 proves the more general statement for every nonnegative input.
The restricted proof above is enough here because the horizon is padded
to at least 16. The ceiling has an exact natural graph:
\begin{equation}\label{eq:direct-ceiling}
 V+B=BF+r+1,\qquad B=r+s+1,\qquad r,s\in\N.
\end{equation}
For $B>0$ these rows give $V+s=BF$ and $0\le s<B$, hence
$F=\lceil V/B\rceil$. If $V\ge1$ they also force $F\ge1$.

\subsection{A radix that can be justified before the factorial}

Define, using natural variables $t_0,x,y,F,b_0$,
\begin{equation}\label{eq:direct-radix}
 \begin{gathered}
 t=t_0+16,\qquad R=2+t+x+y+F,\qquad K=2R,\\
 B=b_0+2,\qquad N=B^K,\qquad q=N^K,\qquad V=(F^2+1)^t.
 \end{gathered}
\end{equation}
These variable powers are explanatory requests: their final implementation
uses the explicit indexed-Pell graph below. Once they have been recovered,
the coordinate shift $B\ge2$ gives
\begin{equation}\label{eq:direct-independent-radix}
 N\ge2^K,\qquad q=B^{K^2}\ge2^{4R^2}.
\end{equation}
No interpretation of $B$ as a binomial is needed for this bound.

For $n=F^2+t$ and $0\le j\le t$,
\begin{equation}\label{eq:direct-first-digit-bound}
 \binom nj\le n^t\le(R^2+R)^t
       \le2^{2Rt}<2^{2R^2}<q.
\end{equation}
Use the zero convention if $j>n$. The elementary bound $R^2+R\le4^R$
holds for $R\ge1$, and $t<R$. Thus the first exact digit graph can be
validated \emph{before} any factorial or CRT meaning is available. It yields
$B=\binom{F^2+t}{t}$; equations~\eqref{eq:direct-ceiling} and
Lemma~\ref{lem:direct-factorial} then yield $F=t!$.

This order removes the potential circularity in identifying the two
binomials. The previous radix $N(N+V+1)^K$ is replaced by $N^K$; its
bootstrap is supplied by an explicit lower bound on $B$, not by assuming
the value of the very digit being tested.

\subsection{Six exact powers and eleven modular Pell residues}

For $a\ge2$, write
\[
 X_0(a)=1,\quad Y_0(a)=0,\quad
 X_{j+1}=aX_j+(a^2-1)Y_j,\quad Y_{j+1}=X_j+aY_j.
\]
The retained compiler supplies indexed-Pell certificates for these pairs,
sharing an auxiliary Pell solution and CRT base across all requested
indices. The final system uses exactly six full power requests:
\begin{equation}\label{eq:direct-six-powers}
 \begin{aligned}
 V&=(F^2+1)^t,& N&=B^K,&q&=N^K,\\
 D_t&=q^t,&D_y&=q^{y+1},&D_x&=q^{x+1}.
 \end{aligned}
\end{equation}
The denominator $D_{t+1}=qD_t$ is obtained by multiplication, so it needs
no separate power request or Pell pair at index $t+1$.

Let $W$ be two plus the sum of the distinct bases, exponents and outputs
in~\eqref{eq:direct-six-powers}. The common frame is
\begin{equation}\label{eq:direct-free-frame}
 L^2-((W+1)^2-1)(WZ)^2=1,\qquad Z\ge1,\qquad A=L+\sigma.
\end{equation}
All coordinates are natural. Pell classification and
$Y_e(W+1)\equiv e\pmod W$ imply $L>W^W$, hence $A>W^W$. Conversely the
index $e=W$ supplies $L,Z$, after which every $A\ge L$ is allowed.

For an exact request $m=n^k$, the compiler asserts
\[
 T_n=2An-n^2-1\ge0,\qquad
 X_k(A)+nY_k(A)=AY_k(A)+m+T_nh,\quad h\ge0.
\]
The first row excludes $n=0$. The frame bounds $m$ and $n^k$ strictly
below $T_n$, and the Pell recurrence gives their congruence modulo $T_n$,
so $m=n^k$. Every exponent is positive by its literal definition.
This is the inherited simultaneous positive-power theorem; separating
the lower bound $L$ from the actual base $A$ permits the congruence below.

Put
\begin{equation}\label{eq:direct-root-modulus}
 b=q+1,\qquad \mathcal M=qD_{t+1}D_xD_y
                           =q^{t+x+y+4},
\end{equation}
and impose one natural-quotient equation
\begin{equation}\label{eq:direct-common-root}
              2A(q+1)=(q+1)^2+1+\mathcal M h_0.
\end{equation}
For every $n\ge0$ it yields
\begin{equation}\label{eq:direct-residue}
 \mathcal R_n=X_n(A)+(q+1-A)Y_n(A)
                    \equiv(q+1)^n\pmod{\mathcal M}.
\end{equation}
Indeed the substitution $s=q+1-A$ respects $s^2=A^2-1$ modulo
$\mathcal M$, and $(A+s)^n=X_n+Y_ns$. Moreover $\mathcal R_n>0$:
for $n>0$, $X_n>(A-1)Y_n$, and for $n=0$ the value is 1.

For every required denominator $d$, both $qd$ and $Nd$ divide
$\mathcal M$: the first follows from its product definition and the
second from $N\mid q$. The one exact binomial uses
\begin{equation}\label{eq:direct-exact-digit}
 \begin{aligned}
 X_n+(q+1)Y_n&=AY_n+r+dB+qd h,\\
 d&=r+s+1,&q&=B+o+1,
 \end{aligned}
\end{equation}
where $n=F^2+t$, $d=D_t$. Every divisibility request
$N\mid\binom nk$ instead uses
\begin{equation}\label{eq:direct-div-digit}
 X_n+(q+1)Y_n=AY_n+r+Nd h,\qquad d=r+s+1,
 \qquad d=q^k.
\end{equation}
The binomial theorem identifies the digit or its divisibility once all
lower coefficients are below $q$. Positivity of $\mathcal R_n$ ensures
that the constructed quotients are nonnegative.

Completeness of the common root is not automatic and must be checked.
As $\mathcal M$ is a power of $q$,
$\gcd(q+1,\mathcal M)=1$. Thus $\gcd(2(q+1),\mathcal M)$ is 1 or 2;
in the second case $(q+1)^2+1$ is even. Equation~\eqref{eq:direct-common-root}
therefore defines a soluble congruence class for $A$, containing arbitrarily
large positive representatives. Choose one above $L$ and large enough to
make $h_0\ge0$. Then supply the standard common-index witnesses at this
base. All semantic values were fixed beforehand.

There are eleven modular requests: one exact binomial and ten ranges.
Their upper indices are positive and need not be included in $W$;
the indexed-Pell graph certifies those indices directly. Fifteen independent
indices suffice for the six exact powers and eleven residues together.

\subsection{The explicit numerical source and its finite proof flow}

For completeness, the arithmetic source is specified here before applying
the inherited calibration. Its constructors are
\[
 \begin{aligned}
 I(a,b)&=2(a+b)(a+b+1)+4b,& M(a,b)&=I(a,b)+1,\\
 A(a,b)&=I(a,b)+2,&E(a,b)&=I(a,b)+3,\\
 N_0(a)&=I(a,1),&X_0(a,b)&=N_0(A(a,N_0(b))).
 \end{aligned}
\]
The subscript here distinguishes the formula constructors from the global
CRT modulus $N$ and Pell abscissas. Define
\[
 \begin{aligned}
 \operatorname{req}(a,k,m)
   &=A(a,I(m,N_0(A(a,I(m,N_0(M(a,k))))))),\\
 s(a,b,k)&=A(k,I(M(k,a),M(k,b))),\\
 u(a,b,k,m)&=N_0(A(a,I(s(a,b,k),\operatorname{req}(b,k,m)))),\\
 \operatorname{clos}(a,b,k,m)
   &=A(b,I(M(b,k),N_0(A(a,I(m,u(a,b,k,m)))))),\\
 v(a,b,c)&=I(M(a,b),A(c,M(a,b))),\\
 h(a,b,c)&=A(c,I(\operatorname{clos}(a,b,c,c),
                          N_0(\operatorname{req}(b,c,c)))),\\
 \operatorname{bundle}(a,b,c)&=N_0(I(v(a,b,c),h(a,b,c))).
 \end{aligned}
\]
The nine conclusion polynomials are
\begin{align*}
 c_1&=X_0(a,E(a,b)),\\
 c_2&=I(A(a,I(b,c)),I(A(a,b),A(a,c))),\\
 c_3&=I(E(a,b),I(M(c,a),M(c,b))),\\
 c_4&=I(A(a,A(b,c)),A(b,A(a,c))),\\
 c_5&=I(E(c,a),I(M(c,b),M(a,b))),\\
 c_6&=I(X_0(a,A(a,b)),b),\\
 c_7&=I(E(c,a),I(E(c,b),A(M(a,b),E(a,b)))),\\
 c_8&=I(I(I(c,d),a),I(I(a,c),I(b,c))),\\
 c_9&=\operatorname{bundle}(a,b,c).
\end{align*}
Their code definitions are matched to the retained C222 calculus.
The first eight are unconditional axiom families. The ninth is admitted
when $c\ne a,b$, represented by
$G=(c-a)^2(c-b)^2-(d+1)^2=0$.

Write $w(z)=1+Fz$. For the eight local columns
$(a,b,c,d,S,S',u,v)$ define
\begin{align*}
 H={}&(S'-S)^2\prod_{j=1}^{8}(S'w(c_j)-S)^2\\
 &\cdot\bigl((S'w(c_9)-S)^2+G^2\bigr)\\
 &\cdot\bigl(S'w(b)-Sw(a)w(I(a,b))\bigr)^2
       \bigl(S'w(A(a,b))-Sw(b)\bigr)^2,\\
 \mathcal L(i)&=H+\bigl(x-S-(1+(i+1)y)u\bigr)^2
                  +\bigl(x-S'-(1+(i+2)y)v\bigr)^2,\\
 J&=\sum_{j=1}^{9}c_j+a+b+I(a,b)+A(a,b)+1.
\end{align*}
The local bounds are $(t,t,t,t,y,y,x,x)$, together with the derived
condition $J\le t$. The endpoint equations are
\[
 x=F+1+(y+1)u_0,\qquad y=F+1+r_0,\qquad
 x=1+(1+(t+1)y)u_t.
\]
Every factor of $H$ is nonnegative. A zero selects truth, a genuine
axiom, modus ponens, or generalization, and gives the appropriate
accumulator identity
\[
 S'w(\text{conclusion})=Sw(\text{premise}_1)w(\text{premise}_2).
\]
Unused premises have truth label 0 and weight one.

Because $F=t!$, the weights $1+jF$ for $1\le j\le t$ are pairwise
coprime: a prime common to weights $j\ne k$ would divide $j-k$, hence
divide $F$, which is impossible. The beta equations identify overlapping
accumulators by uniqueness of remainders bounded by $y$. Their positive
initial value and positive weights propagate positivity. Multiplying the
local identities and cancelling gives
\[
 \prod w(\text{conclusions})=w(1)\prod w(\text{premises}).
\]
Taking valuations at a prime of each positive label's weight gives a
multiplicity excess of one at label 1 and zero at every other positive
label. If 1 were outside the derivable closure, each selected rule with
conclusion outside it would have a premise outside it. Counting these
occurrences contradicts that excess. This argument allows cyclic or
repeated selected rules; it does not assume the supplied flow is a tree.

Conversely expand a finite refutation to a proof tree, choose the horizon
above its length, parameters and all finite $J$-values, and pad with truth
steps. Define $F=t!$ only after this choice. Expanding the pending proof
multiset backwards gives the accumulator products. Choose $y$ large and
divisible by $t!$, then choose $x$ by CRT for the beta moduli. This is the
inherited complete proof-flow construction, with no dependence of the
horizon on the later accumulator sizes.

\subsection{Reciprocal prime-power decoding and the last carry bounds}

With $F=t!$, let
\[
 d_i=1+F^2/i\quad(1\le i\le t),\qquad
 B=\prod_{i=1}^{t}d_i=\binom{F^2+t}{t},\qquad N=B^K.
\]
The $d_i$ are integral, pairwise coprime, and have no prime divisors at
most $t$. These facts follow from $i\mid F$ and
$id_i-jd_j=i-j$. Introduce nine positive encoded representatives,
the eight local columns $A_j$ and index $C$, satisfying
\begin{equation}\label{eq:direct-crt-system}
 \begin{gathered}
 C+\sum_j A_j+g=9N^2,\qquad B\mid C+F^2+1,\\
 N\mid\binom Ct,\qquad N\mid\binom{A_j}{b_j+1},\qquad
 N\mid\mathcal L(C,A),\qquad N\mid\binom{J(A)}{t+1}.
 \end{gathered}
\end{equation}
The $b_j$ are the local bounds. The derived value $J(A)$ is deliberately
excluded from the aggregate bound.

The elementary one-sided range lemma states
\[
             p^e\mid\binom ak\quad\Longrightarrow\quad
                         a\bmod p^e<k.
\]
For $a\ge k$, the $p$-valuation of the binomial is the sum of
$\lfloor a/p^j\rfloor-\lfloor(a-k)/p^j\rfloor-\lfloor k/p^j\rfloor$,
each zero or one. If the interval $[a-k+1,a]$ had no multiple of $p^e$,
all terms from $j=e$ onward would vanish, contradicting valuation at
least $e$. The case $a<k$ is immediate. Its converse is false and is not
used.

Choose a prime $p\mid d_i$, write $m=v_p(d_i)$ and $e=Km$. The index
congruence and $p>t$ identify $C\bmod p^e=i-1$. The other ranges give
bounded residues for the local columns. A retained literal majorant
certificate gives $\mathcal L,J<R^E$, where $E=2^{23}$. Since
$R\ge16!>2^{44}$, elementary induction gives
$R^E<2^R<2^K\le p^e$. Polynomial congruence now yields
$\mathcal L(i-1,\text{local})=0$ and $J(\text{local})\le t$.

For completeness choose the representatives by CRT at precision
$d_i^{2K}$, so they lie in $[1,B^{2K}]=[1,N^2]$. Their falling factorial
numerators have a factor of valuation at least $2Km$; division by $k!$
loses at most $k-1\le K-1$. This proves the required divisibilities by
$p^{Km}$, including the derived polynomial range. Precision $2K$ is
essential to this argument; no false converse to the range lemma is used.

The remaining digit bounds also remain strict for the new radix $q=N^K$.
For $n\le9N^2$, $1\le k\le R-1$, and $j\le k$,
\[
 \binom nj\le n^k\le9^kN^{2k}
      <2^{4R}N^{2k}\le N^{2k+2}\le N^K=q.
\]
For the derived argument the positive-DAG certificate gives
$J(a,b,c,d)\le(n+1)^\ell$ when the four inputs are at most $n$,
with $\ell=2^{22}$. Since $n\le9N^2$ and $N\ge10$, $J(A)\le N^{3\ell}$.
The finite inequality $2\cdot16!\ge3\ell\cdot17$ persists by induction,
so $3\ell(t+1)\le2F<K$. Therefore
\[
 \binom{J(A)}j\le J(A)^{t+1}
                  \le N^{3\ell(t+1)}<N^K=q.
\]
This last strict inequality is necessary: using the old statement
$N^K<q$ would be incorrect for the new radix.

\subsection{Closing the theorem and reproducing the finite object}

Soundness first recovers the six exact powers from the independent Pell
frame. It then uses~\eqref{eq:direct-first-digit-bound} to interpret the
single exact binomial and the factorial ceiling, before invoking any CRT
or proof-label meaning. Only afterward are the ten ranges interpreted,
the bounded local source decoded, and the Horn-flow theorem applied.
The inherited ordinary calibration is
\[
 \mathrm{C222}(1)\ \longleftrightarrow\
 \text{finite ZF refutation}\ \longleftrightarrow\
 \text{finite ZFC refutation}.
\]
Its named interfaces are
\begin{quote}\small\ttfamily
KEndpoint.c222\_iff\_zf\_refutation\\
ConstructibleChoice.zf\_refutation\_iff\_zfc\_refutation.
\end{quote}
The construction uses the numerical-calculus endpoint before the
machine-specific theorem; it does not simulate K215 or instantiate a
Jones universal polynomial.

Completeness chooses, in order, the finite proof and padded horizon;
the factorial, proof-flow and beta data; $B,R,K,N,q$ and the CRT
representatives; the six powers and eleven true binomial residues; and
finally a lower Pell frame and sufficiently large $A$ in the common-root
congruence class. The common-index theorem supplies all auxiliary Pell
witnesses. Every displayed natural quotient is then reconstructed.
The checked natural projection gives a zero of $P_{319}$. Conversely
its restoration certificate returns every deleted coordinate, so it
preserves both existence directions.

The optional square-shell companion has 320 natural variables, 255
quadratic rows, $13\,828$ terms and coefficient height $15\,248\,812$.
It trades one extra variable and more terms for smaller coefficients,
using the retained constructor-preserving syntax transport. The main
319-variable result stays with the original Cantor encoding.

The complete source bundle uses only the Python 3 standard library.
Normal and optimized-interpreter replays rebuild both equations and
compare all twenty generated JSON files byte for byte. The independent
checker verifies every restoration, every final coefficient, all live
coordinates and deliberately corrupted controls. Separate finite
regressions cover $18\,049$ factorial cases, 270 independent radix
bounds, 49 growth budgets, $1\,530$ Pell-residue instances and 516 digit
and divisibility instances. Those checks are arithmetic regressions,
not proofs of unbounded statements or searches for a ZFC refutation.

\clearpage
\section{Eleven integer variables with a short divisibility carrier}
\label{sec:short-carrier}
This section gives the degree-440 equation and two cheaper companions.
Its two essential reductions are to replace a fifth-power carrier by
a first-power carrier, and to replace a variable-dependent index modulus
by $4q^4$. The full supported-modulus interpolation theorem is needed for
the converse. Reusing only the older interpolation modulo the strong
ordinate would leave a gap.

\subsection{The fixed constants and the complete polynomial}
The only unknowns are
\[
 (j,C,z,h,r,v,w,T,c,n,s)\in\Z^{11}.
\]
Use the unchanged source-34 tuple
\[
 \beta,K_0,M_C,M_F,A_c,B_c,
 \qquad \beta=2^{1064921}-1,\quad
 A_c=2^{1113},\quad B_c=2^{666104}.
\]
The fourth port $M_F$ is the \emph{already shifted} mask. The six exact
numerals, their bit lengths, hashes, and executable finite reconstruction
are in \code{norm/coefficients}; their common identity is
\code{mixed-parity-unrounded-1064921}. These are constants, not six extra
unknowns or tunable parameters. The tuple must not be mixed with another
source's smaller coefficients.

Define the following polynomial abbreviations:
\begin{align*}
 J&=6j+1,&q&=\beta J+1,&W&=12w+4,\\
 X&=qW,&Y&=(2v+1)q^3,&a&=(X+1)Y,\\
 H&=4a+3,&\Delta&=(a+2)^2-1,\\
 F&=(K_0+W)C-z(q-1),\\
 R&=\bigl(q(q-F)-C\bigr)(q^2-1)+(M_C+qM_F)J,\\
 k&=R+1+4h q^4,&\eta&=c-Yk,\\
 D_{\mathrm{err}}&=(A_cF-q)^2+(B_cC-q)^2,\\
 \mathcal B&=k^2(q^2-D_{\mathrm{err}})-4q^2\eta^2,\\
 \ell&=XY^2,&D_\tau&=1+\ell(\ell+1)k^2,\\
 D&=3\mathcal B-4,&D_g&=1+4D D_\tau s^2,\\
 P&=D_\tau D_g,\\
 I&=2rc^2,&D_f&=1+\Delta I^2,\\
 Q&=\Delta^2 I^2,&\delta&=Q-1,\\
 V&=D_f(c-R+4IT)-c,&C_y&=QV^2-1,\\
 D_a&=\delta C_y,&K&=X(3c-2X)+2,\\
 z_c&=HT+K,&\mathcal A&=(z_c^2+3H^2)(z_c^3-19H^3),\\
 S&=srH,&Z&=S\bigl(n+(2Q-1)V\bigr)-\mathcal A,\\
 u&=S^2D_f,&b&=4S^2D_a,&e&=16S^2P.
\end{align*}
The equation is
\begin{equation}\label{eq:three-root}
 \boxed{
 \left[(Z^2-u-b-e)^2-4(ub+ue+be)\right]^2
       -64ubeZ^2=0.}
\end{equation}
It contains only fixed nonnegative exponents and integer arithmetic. Its
literal division-free evaluator is
\code{norm/scalar11\_q4\_center.json}: 126 operations and exact degree 440.
The following proof uses radicals to analyze the polynomial; no radicals
or concealed existential witnesses occur in the equation itself.

The constant recipe constructs 943 fixed tiles from the retained literal
430-action table, assigns the fixed bit weights and positional offsets,
and forms the two masks and convolution constant by finite sums. It never
simulates the machine. Its code and complete input data are included; the
check reconstructs all four native ports and separately checks $A_c,B_c$.
The six uncompressed numeral payloads total 615,883 bytes. Thus a short
126-operation evaluator does not mean that the full coefficient
description has only 126 symbols.

\subsection{The algebraic norm and the exceptional coordinates}
Let
\[
 \mathcal F_8(Z;u,b,e)=
 [(Z^2-u-b-e)^2-4(ub+ue+be)]^2-64ubeZ^2.
\]
In a free multiquadratic algebra this is the product of
\[
 Z-\epsilon_1\sqrt u-\epsilon_2\sqrt b-\epsilon_3\sqrt e
 \quad(\epsilon_i\in\{-1,1\}).
\]
This identity holds even when radicands coincide or are negative or zero.
An independent exact checker multiplies all eight factors and compares
all 35 collected coefficients.

The fixed affine front gives $q$ even, $q\equiv2\pmod3$, $|X|\ge64$,
$|Y|\ge4096$, $H$ odd and nonzero, and $\Delta>1$.
Also $D_\tau\ge1$ because $\ell(\ell+1)\ge0$ for every integer $\ell$.
The polynomial
\[
 \pi(x)=(x^2+3)(x^3-19)
\]
has no rational root. It has a root modulo every positive integer, by
the retained elementary prime-power and CRT proof in
\code{norm/source/PRIME\_TWELVE.md}. Since
\[
 \mathcal A=H^5\pi(T+K/H),
\]
we have $\mathcal A\ne0$ for every integer assignment.

If $S=0$, every scaled radicand in~\eqref{eq:three-root} is zero and
the polynomial is $\mathcal A^8\ne0$. Thus $r,s\ne0$ at a zero.
This does \emph{not} yet imply $c\ne0$. The order of recovery matters.
We first prove $D_a>0$ in both possible cases. If $c=0$, then $I=Q=0$
and $D_a=(-1)(-1)=1$. If $c\ne0$, then $|I|\ge2c^2$,
$Q>1$, and $D_f>|c|+1$. As $V=D_f b_0-c$ for an integer $b_0$,
either $b_0=0$ and $|V|=|c|$, or
$|V|\ge D_f-|c|>1$. Thus $C_y,\delta,D_a>0$.

In a vanishing conjugate factor, $P=D_\tau D_g$ is therefore the sole
possibly negative radicand. If it were negative its nonzero imaginary
contribution could not cancel the two real roots and real center.
Since $D_g$ is odd, it follows that $D_g>0$. If $\mathcal B\le1$, then
$D\le-1$ and $D_g\le-3$, a contradiction. Hence $\mathcal B\ge2$.
For $c=0$, however,
\[
 \mathcal B=k^2\{q^2(1-4Y^2)-D_{\mathrm{err}}\}\le0.
\]
Consequently $c\ne0$, and only now do we invoke the strong and auxiliary
Pell equations with $I\ne0$ and $Q>1$.

\subsection{Square recovery and auxiliary integrality}
The normalized vanishing factor is a rational signed sum of
\[
 \sqrt{D_f},\qquad 2\sqrt{D_a},\qquad4\sqrt P.
\]
All three radicands are positive odd integers. In the squarefree
decomposition of their roots, the integer square-root coefficients are
odd. Within a nontrivial squareclass the weights have distinct 2-adic
orders $0,1,2$; the term with least order cannot cancel. Linear independence
of distinct squarefree radicals therefore forces all three roots to be
integers. Furthermore
\[
 \gcd(D_\tau,D_g)=\gcd(D_\tau,1+4D D_\tau s^2)=1,
\]
so $D_\tau=\tau^2$ and $D_g=g^2$ separately.

The square root of $D_a$ initially gives a rational $y$ satisfying
\begin{equation}\label{eq:adjacent-norm}
 QV^2-(Q-1)y^2=1.
\end{equation}
We recall the relevant integrality argument explicitly. In lowest terms
$y=a_1/b_1$, one has $b_1^2\mid Q-1$. For $|V|>1$ the usual
adjacent-coefficient Pell transformation is
\[
 V'=(2Q-1)V-2(Q-1)y,\qquad
 y'=(2Q-1)y-2QV.
\]
After choosing signs, it satisfies $|V'|<|V|$, keeps $V'$ integral, and
preserves the reduced denominator of $y$, since $2Q-1\equiv1\pmod{b_1}$.
The norm excludes $V'=0$. Descent ends at $|V|=1$, where $|y|=1$.
Thus $b_1=1$. All normalized roots in the factor are now integers, and
the factor equation yields $S\mid\mathcal A$.

\subsection{The smaller index modulus}
The retained signed rank proof uses a first-Pell congruence modulo
$|XY|>3R+1$. We may instead use
\[
 M=4q^4>3R+1.
\]
The inequality follows from the unchanged packing bound
$0<R<q^4-2$. Since $4\mid W$,
\[
 2XY^2=2W(2v+1)^2q^7\equiv0\pmod M.
\]
The signed first Pell parameter $P_0=|2XY^2+1|$ is therefore
$\operatorname{sign}(X)\pmod M$. Its classification gives
$|k|=2\psi_{P_0}(n_0)$ and the congruence
\[
 \epsilon\,2n_0\equiv R+1\pmod M,
 \qquad
 \epsilon=\operatorname{sign}(k)\operatorname{sign}(X)^{n_0-1}.
\]
For the positive sign $2n_0\ge R+1$; for the negative sign
$2n_0\ge M-R-1>2R$. The remaining strong/auxiliary folding argument
does not use $h$ separately. It gives a main index $p=R$ with $p>n_0$.
Hence the negative branch is impossible. Both $2n_0$ and $R+1$ lie
strictly between zero and $M$, so
\[
 2n_0=R+1.
\]
Only after this equality is proved do we reduce the same recurrence
modulo $|XY|$. It restores $k\equiv R+1\pmod{XY}$ and the old integral
coordinate $h_{\mathrm{old}}=(k-R-1)/(XY)$. The downstream source
argument is consequently unchanged. Conversely the canonical positive
first-Pell witness has the new congruence modulo $M$, so
$h=(k-R-1)/M$ is integral. This changes neither the witness domain nor
the coefficient tuple.

\subsection{Why fifth-power divisibility suffices}
The strong and auxiliary norms, the actual residues
\[
 V\equiv-c\pmod f,\qquad V\equiv-R\pmod c,
\]
the condition $c^2\mid I$, and the centered budget strip give the
pre-projection conclusions of the retained signed folding theorem:
\[
 R=2t_0+1,\quad t_0\ge24,\quad
 |Y|>\frac{(|X|-3)^{t_0}}2-1.
\]
These conclusions precede the missing main projector; it is not being
assumed in order to prove itself. In particular, with $x=|X|$,
\[
 |H|>x(x-3)^{t_0}>x^{t_0+1}/2^{t_0},\qquad x\ge2^{64}.
\]
From $S\mid\mathcal A$ and $H\mid S$, reduction modulo $H$ gives
$H\mid K^5$. Define $K_q=X(X+2ac)-1$. The exact identity
\[
 K=HXc-2K_q
\]
and oddness of $H$ give $H\mid K_q^5$. The signed Pell projection
congruence is
\[
 2^R K_q\equiv\Pi\pmod H,
 \quad \Pi=(\sigma X-2^R)(2^R\sigma X+1),\quad
 \sigma=\operatorname{sign}(c).
\]
Thus $H\mid\Pi^5$. Each factor of $\Pi$ has absolute value smaller
than $2^{R+1}x$, and hence
\[
 |\Pi|^5<2^{20t_0+20}x^{10}<|H|.
\]
For the last inequality, use
\[
 64(t_0-9)-(21t_0+20)=43t_0-596\ge436.
\]
Therefore $\Pi=0$, its odd second factor is nonzero, and
$\sigma X=2^R$. This recovers the exact old projector and permits the
retained orientation, half-binomial recovery, and source restoration.
The invalid inference $H\mid K^5\Rightarrow H\mid K$ is never used.

\subsection{Completeness with the strong pair fixed}
Start with the canonical outer tuple from a finite refutation. It has
$\mathcal B>64$ and $D_\tau=\tau^2$. Moreover $D_\tau\equiv1\pmod3$,
so $D D_\tau$ is positive and $2\pmod3$, hence a nonsquare. A positive
Pell solution, squared once, supplies $s,g>0$ with
\[
 g^2=1+4D D_\tau s^2.
\]
Choose one positive strong pair with $4sc^2\mid I$ and
$f^2-\Delta I^2=1$, and set $r=I/(2c^2)$. Then $r$ is an integer
divisible by $2s$. Every prime of $sr$ divides $I$, although $sr$ need
not divide $I$.

The supported-modulus interpolation theorem retained from Round 14 says
that, for this \emph{fixed} strong pair, every residue class modulo any
positive modulus whose primes divide $I$ occurs as $T$ in an integral
auxiliary solution~\eqref{eq:adjacent-norm}. Its proof uses a Pell period
$\lambda=4m_0f$ and an index increment
$d=\operatorname{lcm}(\lambda,4I)$. For each $p\mid I$, the resulting
map $h\mapsto T(h)$ satisfies
\[
 v_p(T(h_1)-T(h_2))=v_p(h_1-h_2).
\]
It is therefore a permutation modulo each supported prime power, and
CRT supplies the specified modulus. The full proof, including the
2-adic case and the precise initial index, is retained in
\code{norm/source/POWERED\_SATURATION\_PARENT.md}, Section 5, and
was independently reviewed for this use.

The canonical source already has $H\mid K$; write $\kappa=K/H$.
Choose a root of $\pi(T+\kappa)$ modulo $sr$, and interpolate that
class while leaving $I,f,r$ fixed. Then
\[
 \frac{\mathcal A}{S}
   =\frac{H^4\pi(T+\kappa)}{sr}\in\Z.
\]
Finally choose
\[
 n=\mathcal A/S+f+2\delta y+4\tau g-(2Q-1)V.
\]
One factor of the norm vanishes. The order is outer tuple, guard pair,
fixed strong pair, intersective residue, auxiliary solution, and $n$.
This proves completeness without changing a strong witness after a
residue has been selected.

\subsection{Exact degree and the two cheaper companions}
For the primary equation the important degrees are
\[
\begin{array}{c|rrrrrrrrrr}
\text{register}&k&\mathcal B&D_\tau&P&I&S&D_f&Q&V&D_a\\\hline
\deg&5&20&30&82&3&8&18&30&23&106.
\end{array}
\]
Thus $\deg(u,b,e)=(34,122,98)$. Put
$B_0=Sn-\mathcal A$ and $C_0=S(2Q-1)V$, so $Z=B_0+C_0$.
The universal identity
\[
 b-C_0^2=S^2(-V^2-4Q+4)
\]
gives $\deg(Z^2-b)=96$. Writing $\omega=Z^2-b$, the inner bracket
of~\eqref{eq:three-root} is
\[
 (\omega-u-e)^2-4(ub+ue+be).
\]
Its unique highest-degree contribution is $-4be$, of degree 220.
The correction $64ubeZ^2$ has degree 376. Hence the total degree is
exactly 440, with leading form $16b_0^2e_0^2$. On the authentic
coefficient slice a specific attained coefficient is
\[
 [j^{264}w^{40}v^{52}h^{12}s^{12}r^{24}c^{32}z^4]P_{440}
   =2^{476}3^{306}\beta^{264}>0.
\]
The exponent sum is 440. The checker independently propagates leading
forms through the actual graph and verifies the cancellation identity.
An identical 136-operation graph exposes all cancellations and has
syntactic degree bound 440; the selected 126-operation graph denotes the
same polynomial but has the looser naive syntactic bound 488.

Keeping $k=R+1+hXY$ instead gives degree 464 and 125 operations. Removing
the center shear as well gives degree 488 and 121 operations. Translation
of the unrestricted coordinate $n$ preserves zero existence exactly.
The two auxiliary identical graphs of degrees 464 and 440 are audit
artifacts, not additional frontier claims.

One charged operation is saved by the exact identity
\[
 P=D_\tau+D(2D_\tau s)^2
   =D_\tau(1+4D D_\tau s^2).
\]
The short-modulus evaluator also shares $4q^2$ with the budget, computing
$4q^4=(4q^2)q^2$ with one additional multiplication.
The actual operation counts for the three selected graphs are
\[
\begin{array}{c|rrr|r}
\text{degree}&\text{multiplications}&\text{additions}&\text{subtractions}&\text{total}\\\hline
488&77&23&21&121\\
464&79&24&22&125\\
440&80&24&22&126.
\end{array}
\]

\clearpage
\section{A single product equation with 111 arithmetic operations}
\label{sec:product-equation}
The three-root norm can be replaced by one equality between products.
The price is a higher total degree. The resulting two equations have
eleven unrestricted integer variables, using the same entire source-34
coefficient tuple as Section~\ref{sec:short-carrier}:
\[
\begin{array}{l|rr}
\text{first-index modulus}&\text{degree}&\text{operations}\\\hline
XY&966&111\\
4q^4&882&112.
\end{array}
\]
Both improve the inspected all-root fusion equation of degree 1632 in
114 operations. The degree-440 construction remains the lower-degree
alternative at a greater arithmetic cost.

\subsection{Definition of the two polynomials}
Keep the eleven inputs $(j,C,z,h,r,v,w,T,c,n,s)\in\Z^{11}$ and the
common definitions through $P=D_\tau D_g$ in
Section~\ref{sec:short-carrier}. For the 111-operation equation use
$k=R+1+hXY$; for the 112-operation equation use $k=R+1+4hq^4$.
Thus, in both cases,
\[
 D_\tau=1+XY^2(XY^2+1)k^2,
 \qquad D_g=1+4(3\mathcal B-4)D_\tau s^2.
\]
Replace the strong and auxiliary definitions by
\begin{align*}
 A&=a+2,&D_m&=1+\Delta c^2,&U&=rc^2 A D_m,\\
 D_f&=1+\Delta U^2P,&Q&=\Delta^2U^2P,&\delta&=Q-1,\\
 b_0&=c-R+4UPT,&V&=D_fb_0-c,&C_y&=QV^2-1.
\end{align*}
Use the same short projector and scalar carrier,
\begin{align*}
 K&=X(3c-2X)+2,&S&=srH,&z_c&=HT+K,\\
 \mathcal A&=(z_c^2+3H^2)(z_c^3-19H^3),&&
 Z=Sn-\mathcal A.
\end{align*}
The requested equation is simply
\begin{equation}\label{eq:single-product}
 \boxed{\delta Z^2-S^2P D_f C_y=0.}
\end{equation}
All abbreviations are polynomials and introduce no extra variables.
In the literal evaluator,
\[
 P=D_\tau+(3\mathcal B-4)(2D_\tau s)^2,
\]
The short-index graph reuses $4q^2$ from the budget to form
$(4q^2)q^2=4q^4$, saving one more multiplication. The product $UP$ is shared between
$\Delta U^2P=(\Delta U)(UP)$ and $b_0=c-R+4(UP)T$.
The complete graphs are \code{lowops/product11\_degree966.json} and
\code{lowops/product11\_degree882.json}. Each is bound to all six original
fixed numerals; no coefficient parameter remains free.

\begin{theorem}\label{thm:single-product}
Each polynomial in~\eqref{eq:single-product} has an integer zero exactly
when the retained native source has a finite refutation. Composing its
ordinary source calibration gives
\[
 \ZFC\vdash
 \bigl(\text{no zero of the specified polynomial in }\Z^{11}\bigr)
 \longleftrightarrow\Con(\ZFC).
\]
The exact degrees and arithmetic counts are the two rows displayed above.
\end{theorem}

\subsection{Coprimality recovers the signs as well as the squares}
The unconditional affine bounds from the preceding section give
$A\ne0$, $H\ne0$ odd, $\Delta=A^2-1>1$, and $X\equiv0\pmod8$.
For every integer $\ell$, $\ell(\ell+1)\ge0$. Therefore
\[
 D_\tau\ge1,\quad D_\tau\equiv1\pmod8,\quad
 D_g\equiv1\pmod4,\quad P\equiv1\pmod4.
\]
In particular $P\ne0$. Neither $D_f$ nor $\delta$ can vanish, because
$\Delta>1$. Also $C_y\ne0$: the integer identity $QV^2=1$ would force
$Q=1$, which is impossible. Put $D_a=\delta C_y$; it too is nonzero.
These statements do not assume that $P,Q$, or $D_f$ are positive.

If $S=0$, the polynomial becomes $\delta\mathcal A^2$, which is nonzero.
Indeed
\[
 \mathcal A=H^5\pi(T+K/H),\qquad
 \pi(t)=(t^2+3)(t^3-19),
\]
and $\pi$ has no rational root. Thus a zero has $S\ne0$, so $s,r\ne0$.
We deliberately leave $c=0$ open until the guard is recovered.

The three integers $P,D_f,D_a$ are pairwise coprime. The first two
congruence computations are immediate:
\[
 D_f\equiv1\pmod P,\qquad D_a\equiv1\pmod P.
\]
For the third, reduction modulo $D_f$ gives
\[
 Q\equiv-\Delta,\quad\delta\equiv-A^2,\quad
 V\equiv-c,\quad C_y\equiv-D_m,\quad D_a\equiv A^2D_m.
\]
But $AD_m\mid U$, so $D_f\equiv1$ modulo both $A$ and $D_m$.
Consequently $\gcd(D_f,D_a)=1$. All gcds refer to absolute values;
the computations hold at negative inputs and at $U=0$.

Multiply~\eqref{eq:single-product} by $\delta$ and divide by $S^2$:
\begin{equation}\label{eq:coprime-square-product}
 \left(\frac{\delta Z}{S}\right)^2=P D_fD_a.
\end{equation}
A rational square root of an integer is integral. Pairwise coprimality
therefore makes $|P|,|D_f|,|D_a|$ separately squares. Since
$P\equiv1\pmod4$, it cannot be a negative odd square. Hence $P>0$,
then $D_f=1+\Delta U^2P>0$, and then~\eqref{eq:coprime-square-product}
forces $D_a>0$. There are positive integers
\[
 p_0^2=P,\qquad f^2=D_f,\qquad e_0^2=D_a.
\]
The further identity $\gcd(D_\tau,D_g)=1$ splits $P$ once more:
\[
 D_\tau=\tau^2,\qquad D_g=g^2,\qquad p_0=\tau g.
\]
Because $s\ne0$, the inequality $\mathcal B\le1$ would imply
$D_g\le1-4D_\tau<0$. Thus $\mathcal B\ge2$. The centered budget now
gives the common source signs and $|c/k-Y|<1/2$. It also excludes $c=0$:
at that value,
\[
 \mathcal B=k^2\bigl(q^2-D_{\rm err}-4q^2Y^2\bigr)\le0.
\]
We have obtained $c\ne0$ and $U\ne0$ without an initial assumption about
the signs of the saturated radicands.

\subsection{Recovering the original Pell witnesses}
Define, only within the proof,
\[
 I=Up_0,\qquad T_{\rm old}=p_0T.
\]
These are integers, $I\ne0$, $c^2\mid I$, and $|I|\ge c^2$.
The restored equalities are exact:
\begin{align*}
 f^2&=1+\Delta I^2,&Q&=(\Delta I)^2>1,\\
 4IT_{\rm old}&=4UPT,&
 V&=f^2(c-R+4IT_{\rm old})-c.
\end{align*}
Moreover $y=e_0/\delta\in\Q$ satisfies
$QV^2-(Q-1)y^2=1$. The adjacent-coefficient descent used in
Section~\ref{sec:short-carrier} makes $y$ integral. Its hypotheses
$Q>1$ and integer $V$ have now been established. Cancellation of the
nonzero $\delta$ in~\eqref{eq:single-product} yields
\[
 Z^2=S^2p_0^2f^2y^2,
\]
so $Z/S$ is an integer. Hence $S\mid\mathcal A$ and $H\mid K^5$.

The complete signed folding and rank argument now has its original
hypotheses: the budget strip, the first Pell square, a nonzero strong
ordinate divisible by $c^2$, and the auxiliary norm with
$V\equiv-c\pmod f$, $V\equiv-R\pmod c$. Its use does not require
an odd strong multiplier. In the $4q^4$ version the first-index argument
is precisely the one already proved in Section~\ref{sec:short-carrier}.
The exact index and orientation later restore
$XY\mid k-R-1$, so the original quotient $h$ can be regenerated.

The fifth-power height argument also applies unchanged. With
$R=2t+1$, $t\ge24$, $x=|X|\ge2^{64}$ and
\[
 \Pi=(\operatorname{sign}(c)X-2^R)
          (2^R\operatorname{sign}(c)X+1),
\]
it gives $H\mid\Pi^5$ and $|\Pi|^5<|H|$. Thus $\Pi=0$, the exact
main projection is restored, and the signed half-binomial argument
recovers every outer native condition. Fresh private witnesses supply
an actual zero of the retained source. This proves the forward
finite-refutation implication using the same source calibration.

\subsection{Completeness and the extra interpolation factor}
Start with a canonical outer source witness. It has $A>1$ even,
$R\equiv3\pmod4$, $c=\psi_A(R)>0$, $c\equiv R\pmod4$,
$D_\tau=\tau^2$, and the stronger budget $\mathcal B\ge64$.
The positive nonsquare $(3\mathcal B-4)D_\tau\equiv2\pmod3$
has a Pell unit. Squaring one supplies positive $s,g$ with
$g^2=1+4(3\mathcal B-4)D_\tau s^2$. Fix them and put $p_0=\tau g$.

Choose one positive strong pair $I,f$ satisfying
\[
 4sAD_m p_0c^2\mid I,
 \qquad f^2-\Delta I^2=1.
\]
Set $r=I/(AD_mp_0c^2)$ and $U=I/p_0$. The strong pair is now fixed.
Every prime factor of $M=p_0sr$ divides $I$. The canonical projector
gives $\kappa=K/H\in\Z$. Choose an intersective residue $a_0$ modulo
$sr$ with $\pi(a_0+\kappa)\equiv0\pmod{sr}$. Supported-modulus
interpolation, at this same fixed $I,f$, supplies the full auxiliary
witness with
\[
 T_{\rm old}\equiv p_0a_0\pmod{p_0sr}.
\]
The extra factor $p_0$ in this modulus is essential. It makes
$T=T_{\rm old}/p_0$ an integer and ensures $T\equiv a_0\pmod{sr}$.
Therefore
\[
 \frac{\mathcal A}{S}
     =\frac{H^4\pi(T+\kappa)}{sr}\in\Z.
\]
Finally set $n=\mathcal A/S+p_0fy$. Equation~\eqref{eq:single-product}
vanishes. The order of choices is outer witness, $s,g$, one strong
$I,f,r$, a residue, $T_{\rm old},T,y$, and $n$. No coprimality between
these moduli is assumed, and the strong pair is never changed after
the residue request. This proves the reverse implication.

\subsection{Exact degree and finite validation}
The key degrees in the two literal graphs are
\[
\begin{array}{c|rrrrrrrrrr}
 &k&P&U&D_f&Q&b_0&V&C_y&S&\mathcal A\\\hline
XY&7&94&23&152&164&118&270&704&8&35\\
4q^4&5&82&23&140&152&106&246&644&8&35.
\end{array}
\]
In both cases $\deg Z=35$. The summand $\delta Z^2$ has degree 234
or 222, while $S^2PD_fC_y$ has degree 966 or 882. The latter therefore
determines the full leading homogeneous part.

Put $q_0=6\beta j$, $X_0=12q_0w$, $Y_0=2vq_0^3$,
$a_0=X_0Y_0$, and
\[
 k_0=ha_0\quad\hbox{or}\quad
 k_0=4hq_0^4-q_0^3(12wC-q_0z),
\]
respectively. Direct exact multiplication gives the leading form
\[
 2^{36}3^7q_0^{14}X_0^{70}Y_0^{112}k_0^{42}
       s^{16}r^{12}c^{40}T^2.
\]
For the native modulus its coefficient at
$j^{588}w^{112}v^{154}h^{42}s^{16}r^{12}c^{40}T^2$ is
$2^{1002}3^{707}\beta^{588}>0$. For the short modulus the coefficient
at $j^{588}w^{70}v^{112}h^{42}s^{16}r^{12}c^{40}T^2$ is
$2^{960}3^{665}\beta^{588}>0$. Their exponent sums give the two
stated exact degrees. The latter coefficient is selected by the maximal
$h$ power, so no cancellation from the other leading summands is possible.

The independent checker uses only the Python standard library and
does not import the producer. It verifies every local polynomial formula,
all 135 or 136 live nodes, all six complete coefficient payloads,
the coprimality remainder identities, both saturation identities,
the new index-modulus identity, and exact leading monomials at every
gate. Five altered mathematical formulas in each graph are rejected.
Normal and optimized Python runs agree, and both graphs rebuild
byte-for-byte. These finite checks certify the exported algebra and
counts; the unbounded semantic argument is the proof above and its
explicitly retained source lemmas.

\clearpage
\section{Reproduction, validation, and remaining work}
\label{sec:validation}

\subsection{What the archive contains}
The root \code{README.md} is the execution guide. The compiled PDF and
all its editable TeX sections are included. The mathematical artifacts
are organized as follows.
\begin{description}[style=nextline]
\item[\code{direct/}]
The literal 319-variable quartic, all 254 quadratic residuals, its fully
collected 10,283-term coefficient list, readable versions of both lists,
the raw source, and every natural-valued restoration polynomial. The
builder, independent verifier, source theorem, imported compiler code,
original notices, source ledger, and comparison receipt are supplied.
The optional \code{equation320/} companion uses 320 naturals and 255
residuals, has 13,828 terms and height 15,248,812; it trades one variable
and more terms for much smaller height through the inherited square-shell
syntax translation.
\item[\code{norm/}]
The three selected eleven-integer DAGs of degrees 488, 464, and 440;
two polynomial-identical cancellation-exposing audit DAGs; the six
literal coefficient numerals and their reconstruction recipe; all
retained ordinary proof dependencies; and the new exact finite checker.
\item[\code{lowops/}]
The two product-form eleven-integer equations of degrees 966 and 882,
their ordinary proof, literal graphs, independent polynomial and degree
audit, and source bindings to the same fixed tuple.
\item[\code{reviews/}]
Separate mathematical reviews of the three new transformations. These
are independent research-agent audits within this workflow, not an
external human endorsement. Their scope is stated in each file.
\end{description}
The top-level SHA-256 manifest covers the shipped files. The hashes are
integrity checks, not substitutes for the actual numeral data or proofs.

\subsection{The direct finite audit}
The principal raw system has 465 natural coordinates and 400 quadratic
rows. The delivered contraction removes 146 coordinates and 146 rows.
For each deletion the certificate gives a polynomial restoration that is
nonnegative on the remaining natural coordinates. The independent checker
restores each deleted row exactly, verifies every substitution, expands
the entire sum of squares, and compares every collected coefficient.
This proves a finite equation-to-equation equivalence independently of
the number-theoretic interpretation of the raw system.

The readable residual and polynomial files are parsed back into the
same coefficients. Altered residual, restoration, and coefficient controls
are rejected. Both normal and optimized Python execution rebuild the
principal and companion artifacts, comparing all 20 generated JSON
files byte for byte in temporary directories. The complete 349-variable
baseline is also re-expanded and its contraction replayed, so the table
does not rely only on an inherited headline count.

The semantic regression suite checks 18,049 factorial cases, 270 radix
bounds, 49 growth budgets, 1,530 modular Pell residue instances, and 516
digit/divisibility instances. These finite cases are useful implementation
checks. The unbounded lemmas are proved in the report and retained source
proofs; their validity is not inferred from these sample counts.

\subsection{The small-variable finite audits}
The norm checker does not import its producer. It binds every live
register to an independently transcribed formula in a free polynomial
ring, checks all 35 coefficients of the three-root norm, proves the
center cancellation and short-projector identities, checks topology and
liveness, and propagates an exact nonzero leading coefficient through
each actual graph. The primary degree-440 coefficient displayed in
Section~\ref{sec:short-carrier} is therefore an attained coefficient on
the actual fixed-coefficient slice.

The product-form checker likewise independently binds the complete
register definitions, verifies the pairwise-coprimality congruences and
the exact leading induction, and counts every live operation. A separate
review rederived its leading forms and confirmed the final
111/112-operation census. Reusing the existing $4q^2$ register saves
one operation in the second graph without changing its polynomial.

All new checks use explicit exceptions rather than removable Python
assertions. The inherited coefficient constructor deliberately requires
normal execution; the normal audit rebuilds it and compares every native
coefficient numerically with the authenticated payload. Optimized checks
retain the same payload authentication and all new exact algebra checks.
The archive does not bypass that constructor's execution contract.

\subsection{Offline commands}
Only Python 3's standard library is needed for the polynomial artifacts.
The root reproduction wrapper supplies the exact command set; the main
component commands are:
\begin{lstlisting}
python3 -B direct/verify.py
python3 -B direct/check_semantics.py
python3 -B direct/reproduce.py
python3 -B -O direct/reproduce.py
python3 -B norm/check.py
python3 -B -O norm/check.py
python3 -B lowops/check.py
python3 -B -O lowops/check.py
\end{lstlisting}
No command searches for a root or decides consistency. The TeX source can
be compiled with a standard TeX Live installation, running
\code{pdflatex -interaction=nonstopmode -halt-on-error report.tex}
in \code{report/} until references stabilize.

\subsection{What remains open}
The strongest unresolved issue is a genuine lower bound or a new
architecture that improves the small-variable and small-coefficient
objectives together. This work does not reduce the minimum variable
count already achieved elsewhere in the repository. It improves the
eleven-integer degree/circuit tradeoffs and a direct modest-height
quartic architecture.

Two next steps are concrete. First, internalize the new signed product
recovery and supported-modulus witness construction in the proof
assistant, keeping its calibration theorem explicit. Second, investigate
whether the shared factorial/CRT certificate admits a further reduction
of the independent index certificates or bounded local columns. Any such
claim must preserve the noncircular digit proof and the complete
refutation witness order. Merely hiding a certificate inside an
unexpanded abbreviation does not reduce its charged complexity.



\clearpage
\begin{thebibliography}{9}
\raggedright
\bibitem{proveit} V. Reshetnikov and ProveIt contributors,
\emph{ProveIt}, pinned commit
\texttt{45b399ec51dfbecb891dd51d8a9e4b6b7c361ac6}, 24 October 2026.
\url{https://github.com/VladimirReshetnikov/ProveIt}.
The precise parent proof texts, numerical payloads, and source notices
used here are retained in the archive.
\bibitem{incoming} ProveIt incoming research packets,
\emph{Coprime Projection}, \emph{Powered Saturation and Factorial Carriers,
Round 14}, \emph{Synchronized Exponents}, and the other incoming packets
listed in the provenance inventory, 24 October 2026.
\url{https://github.com/VladimirReshetnikov/ProveIt/tree/45b399ec51dfbecb891dd51d8a9e4b6b7c361ac6/docs/incoming}.
\bibitem{jones} J. P. Jones, Universal Diophantine equation,
\emph{Journal of Symbolic Logic} \textbf{47} (1982), 549--571.
\url{https://doi.org/10.2307/2273588}.
\bibitem{sun} Z.-W. Sun, Further results on Hilbert's tenth problem,
\emph{Science China Mathematics} \textbf{64} (2021), 281--306.
\url{https://arxiv.org/abs/1704.03504}.
\bibitem{bayer} J. Bayer, M. David, M. Hassler, Yu. Matiyasevich,
and D. Schleicher, Diophantine equations over $\Z$: universal bounds and
parallel formalization, arXiv:2506.20909.
\url{https://arxiv.org/abs/2506.20909}.
\bibitem{afp} M. David and collaborators, Universal Pairs for Diophantine
Equations, \emph{Archive of Formal Proofs}, 2025.
\url{https://isa-afp.org/entries/Diophantine_Universal_Pairs.html}.
\end{thebibliography}
\end{document}
